\section{Preliminaries}

\subsection{Symbolic Planning}
\subsubsection{Classical Planning with Types}

% We adopt the PDDL planning formalism~\cite{mcdermott1998pddl}, which extends classical STRIPS~\cite{fikes1971strips} with a type hierarchy. A planning domain $D$ consists of a set of predicates $\mathcal{P}$, a set of actions $\mathcal{A}$, and a type hierarchy $\mathcal{T}$. Each type $\tau \in \mathcal{T}$ denotes a set of objects; the hierarchy imposes $\tau_1 \sqsubset \tau_2$ (e.g., $\texttt{eraser} \sqsubset \texttt{tool} \sqsubset \texttt{object}$). Each action $a \in \mathcal{A}$ is defined as $a = \langle \text{pre}(a), \text{add}(a), \text{del}(a) \rangle$ where $\text{pre}(a), \text{add}(a), \text{del}(a) \subseteq \mathcal{P}$, and each parameter of $a$ is constrained to a type $\tau \in \mathcal{T}$. A planning problem $\Pi = (D, O, \sigma, I, G)$ pairs the domain with a set of objects $O$ and a type assignment $\sigma : O \rightarrow \mathcal{T}$. Let $\mathcal{F}$ denote the set of ground atoms obtained by instantiating the predicates $\mathcal{P}$ with objects from $O$. A \emph{state} $s \subseteq \mathcal{F}$ is a set of ground atoms; the initial state $I \subseteq \mathcal{F}$ is one such state, and the goal $G \subseteq \mathcal{F}$ is a set of atoms defining the \emph{goal states} $\{\, s \mid G \subseteq s \,\}$. The type assignment $\sigma$ determines action applicability: an action schema parameterised over type $\tau$ may only be grounded with objects $o$ such that $\sigma(o) \sqsubseteq \tau$.

% A \textit{plan} $\pi$ is a sequence of grounded actions $\langle a_1, \ldots, a_k \rangle$ such that executing the sequence from $I$ satisfies $G$. A symbolic planner takes $\pi$ as input and returns either a valid plan or reports that no plan exists.


We adopt the typed STRIPS fragment of the Planning Domain Definition Language (PDDL)~\cite{mcdermott1998pddl}, which extends classical STRIPS~\cite{fikes1971strips} with typed objects and action parameters. A planning domain $D = \langle \mathcal{P}, \mathcal{A}, \mathcal{T}, \sqsubseteq \rangle$ consists of predicate symbols $\mathcal{P}$, action schemas $\mathcal{A}$, types $\mathcal{T}$, and a reflexive subtype relation $\sqsubseteq$ (e.g., $\texttt{eraser} \sqsubseteq \texttt{tool} \sqsubseteq \texttt{object}$). Each action schema $a \in \mathcal{A}$ is a tuple $a = \langle \mathrm{par}(a), \mathrm{pre}(a), \mathrm{add}(a), \mathrm{del}(a) \rangle$, where $\mathrm{par}(a)$ are typed parameters and the remaining components are sets of atoms over those parameters. Preconditions are positive conjunctions, and effects are represented by add and delete lists.

A planning problem $\Pi = \langle D, O, \sigma, I, G \rangle$ consists of a domain $D$, a finite set of objects $O$, a type assignment $\sigma : O \rightarrow \mathcal{T}$, an initial state $I$, and a goal $G$. Let $\mathcal{F}$ be the set of all well-typed ground atoms over $\mathcal{P}$ and $O$. A state $s \subseteq \mathcal{F}$ contains the atoms true in that state, while omitted atoms are false under the closed-world assumption. The initial state $I \subseteq \mathcal{F}$ is a single fully specified state, whereas $G \subseteq \mathcal{F}$ defines the goal-state set $S_G = \{\,s \subseteq \mathcal{F} \mid G \subseteq s\,\}$.

A parameter of type $\tau$ may be grounded with an object $o$ only if $\sigma(o) \sqsubseteq \tau$. A well-typed grounded action $a$ is applicable in $s$ if $\mathrm{pre}(a) \subseteq s$, and produces the successor state $\gamma(s,a) = (s \setminus \mathrm{del}(a)) \cup \mathrm{add}(a)$. A plan for $\Pi$ is a finite sequence of well-typed grounded actions $\rho = \langle a_1,\ldots,a_k\rangle$ that are successively applicable and satisfy $G \subseteq \gamma(I,\rho)$. A sound and complete symbolic planner takes $\Pi$ as input and returns a valid plan if one exists, or establishes that no plan exists.

\label{sec:preliminaries}

\subsection{Running Example}
Consider the task depicted in Figure~\ref{fig:robot_example}. The environment consists objects $\mathcal{O} = \{\mathtt{marker}, \mathtt{cloth}, \mathtt{cup}, \mathtt{blackboard}\}$. The robot's task is to wipe the blackboard, and it has actions $\mathcal{A} = \{\mathtt{wipe}, \mathtt{pickup} \}$. 

% We also assume the same formalisms defined in the Symbolic Planning section. Therefore, 
The actions and objects possess types which restricts the number of grounded actions generated for each lifted action encoding. For example, an object can be typed as \texttt{pickable}, which simply categorizes objects that can be assigned to the \texttt{pickup} action. Let us assume that these types correspond to the object names listed in $\mathcal{O}$. In the example, no object of type \texttt{eraser} is present, so a plan would not be able to be generated. However, in this example a different object, \texttt{cloth} could be used in place of an \texttt{eraser}. 
 

% For example, the \texttt{wipe} action can encode that an object of the type \texttt{eraser} can be used to erase an object of the type \texttt{blackboard}. Each object would also be assigned a type, let us assume that these types correspond to the object names listed in $\mathcal{O}$.

% We consider two ways the the set of object can be generated:  they may be manually added to the problem file, or the problem file might be generated at run time using the robots sensors such as camera. In the case where a hand crafted problem file is available, objects must be annotated with labels that allow the planner to determine which actions can be applied to which objects. This requires a comprehensive list of types, and additionally may not encode different ways of solving the problem apriori. In the other scenario where the problem file is generated from sensor observations - it is important to consider how types can be associated with each object. In the most general sense, a problem file can be generated assigning each object to a generic type, however this would allow for plans like using the `cup` to clean the `blackboard`, which is not reasonable.


\subsubsection{The Typing Problem}

A \textit{typed domain} $D_t$ encodes affordances via type constraints: for example, \texttt{wipe} requires an $\texttt{eraser}$-typed tool and a $\texttt{whiteboard}$-typed surface. In practice, objects such as \texttt{cloth} may be assigned only the base type $\texttt{tool}$ --- either because the problem was constructed at runtime from untyped perception, or because a fixed object-to-type mapping is unavailable. An object typed at $\texttt{tool}$ is ineligible for any action whose parameter requires $\texttt{eraser}$, even if that object could functionally serve the role. Not every object in $O$ requires reassignment --- objects already typed at an appropriate subtype, or irrelevant to the goal's action chain, retain their types under $\sigma$.

Crucially, the goal $G$ determines which types are required; for example to achieve $\texttt{(clean blackboard)}$, the planner must execute $\texttt{wipe}$, which requires an $\texttt{eraser}$-typed object. Observe that properties irrelevant to this action chain, such as an objects $\texttt{color}$ ,need not be typed. When $\sigma$ assigns objects to insufficient types, the grounded action space is such that no plan achieving $G$ exists. We call this a \textit{typing gap}. A typing gap is a restricted instance of \emph{model repair}~\cite{bercher2025modelrepair}: the only modification permitted is reassignment of object types, leaving the domain, goal, and initial state untouched.

\subsubsection{Failure Modes}

Given a typed domain $D_t$ and a problem instance with a typing gap, there are two distinct failure cases. First, a valid substitution may exist --- some object in $O$ could plausibly fill the required subtype --- but the type assignment $\sigma$ fails to reflect this. Second, no object in $O$ is a plausible instance of the required subtype: no valid substitution exists and the correct behavior is to return no plan.

\subsection{Large Language Model Queries}
\label{subsec:llm_query}

Given a prompt $\rho$ encoding a planning domain $D$ and problem $\pi$, a large language model $L$ returns a proposed type assignment $\hat{\sigma} : O \rightarrow \mathcal{T}$, expressed as a set of pairs $\{(o, \tau) \mid o \in O,\, \tau \in \mathcal{T}\}$. Since $L$ is a stochastic function, repeated queries over the same $\rho$ may yield different assignments. The goal of LLM-based typing is to find a $\hat{\sigma}$ that resolves the typing gap in $\pi$, enabling a symbolic planner to return a valid plan.

\subsection{Problem Formulation}

Given a problem $\pi$ exhibiting a typing gap, our objective is to find a type assignment $\hat{\sigma}$ under which the planner returns a plan achieving $G$, or to correctly report that no valid substitution exists. Every strategy in the next section instantiates the same template: an LLM proposes one or more candidate assignments $\hat{\sigma}$, and each is verified by invoking the symbolic planner on $\pi|_{\hat{\sigma}}$. Because verification is delegated to the planner rather than the LLM, the framework inherits the planner's guarantees relative to the assigned types.

While we retain the planner's formal guarantees, this is not to say the approach is without a trade-off. In exchange for the flexibility to solve problems a rigid planner would declare unsolvable, we give up any guarantee on the \emph{semantic} validity of $\hat{\sigma}$ --- whether a substitution such as \texttt{cloth}-as-\texttt{eraser} is physically warranted. The planner cannot certify this; it is precisely what our experiments serve to measure.